% Copyright (c) 2026 Eugene Joseph M. Ragasa.
% SPDX-License-Identifier: Apache-2.0
% Author: Eugene Joseph M. Ragasa
% Formal proof completed: 2026-10-19
\documentclass[11pt]{article}
\usepackage[T1]{fontenc}\usepackage{amsmath,amssymb}\usepackage[margin=1in]{geometry}\usepackage[hidelinks]{hyperref}
\title{Proof Dossier: Ragasa2026A Canonical Loss-Set Distance}\author{Eugene Joseph M. Ragasa}\date{Formal proof completed: 2026-10-19}
\begin{document}\maketitle
\section{Identity and status}
ID \texttt{BG-RAGASA2026A-THM-024}; status \texttt{proved}. Lean declaration:
\par\noindent {\tiny\path{ProjectKoios.BandGauge.Papers.Ragasa2026A.}}\\
{\tiny\path{canonicalLossAdmissibleSetDist_eq}}.
\section{Exact statement}
For nonempty canonical loss sets, their Euclidean set distance is exactly the existing
closed-form full-ellipse distance with center $\mu_{\max}$ and the derived operator
radius.
\section{Source, proof, and boundary}
This \texttt{formalizes} the full-plane distance clause near manuscript line 737 at owner
commit
\begin{center}\small\texttt{99a46cb59d0c6f27ada7d709f194996ef9619d9f}.\end{center}
The proof transports the reusable exact-distance theorem across the two field-to-scalar
loss equalities. Set nonemptiness is explicit; no numerical threshold is certified.
Kernel axioms are \texttt{propext}, \texttt{Classical.choice}, and
\texttt{Quot.sound}; no admissions.
\end{document}
